\begin{figure}[H]
\centering
\begin{tikzpicture}[
  fase/.style={
    rectangle,
    rounded corners=3pt,
    draw=AzulUSTA,
    line width=0.7pt,
    fill=AzulUSTA!8,
    text width=0.72\textwidth,
    minimum height=1.25cm,
    align=center,
    inner sep=7pt
  },
  flecha/.style={-{Stealth[length=3mm,width=2mm]}, line width=0.8pt, draw=AzulUSTA}
]
  \node[fase] (identificacion) {\textbf{1. Identificación}\\
  \node[fase, below=8mm of identificacion] (analisis) {\textbf{2. Análisis}\\
  \node[fase, below=8mm of analisis] (diseno) {\textbf{3. Diseño}\\
  \node[fase, below=8mm of diseno] (implementacion) {\textbf{4. Implementación}\\
  \node[fase, below=8mm of implementacion] (validacion) {\textbf{5. Validación}\\
  \node[fase, below=8mm of validacion] (seguimiento) {\textbf{6. Seguimiento y mejora}\\

  \draw[flecha] (identificacion) -- (analisis);
  \draw[flecha] (analisis) -- (diseno);
  \draw[flecha] (diseno) -- (implementacion);
  \draw[flecha] (implementacion) -- (validacion);
  \draw[flecha] (validacion) -- (seguimiento);
  \draw[flecha] (seguimiento.east) to[out=0,in=0,looseness=1.2]
    node[right, align=left, font=\footnotesize\itshape, text=AzulUSTA]
    {Retroalimentación para la mejora continua} (identificacion.east);
\end{tikzpicture}
\caption{Secuencia general del ciclo de mejora continua aplicado al proceso CIP.}
\label{fig:secuencia-cip}
\end{figure}